\section{Overview} \label{sec:overview}

The \texttt{smc2fc} package implements two coupled algorithms: a
Sequential Monte Carlo$^2$ (SMC$^2$) Bayesian filter for nonlinear,
non-Gaussian state-space models with intractable marginal likelihoods,
and a stochastic-optimal-control receding-horizon controller built on
the same machinery. The two share the same outer-loop tempered SMC and
the same per-particle differentiable likelihood evaluator, applied to
different log-densities: the Bayes posterior in the filter case, the
control-as-inference posterior in the controller case.

This document describes what is implemented today, with enough
mathematical context to read the source files without an external
textbook open in parallel.

\subsection{Scope and non-scope}

\paragraph{In scope.} The framework code under
\path{smc2fc/}: the core SMC$^2$ engine (\path{smc2fc/core/}), the inner
particle filter (\path{smc2fc/filtering/}), the control-spec abstraction
and the tempered-control loop (\path{smc2fc/control/}). All
mathematical content here applies to any user model that conforms to
the framework's two interfaces (Section~\ref{sec:setting}).

\paragraph{Out of scope.} Specific user models. \texttt{smc2fc} is
designed to be coupled to a user-supplied SDE plus observation model;
the FSA family of physiological models in the sibling
\texttt{FSA\_model\_dev} repository is one such model and is documented
separately. Nothing in this document depends on the FSA model's drift,
state space, or parameters.

\subsection{Architecture in one diagram} \label{sec:architecture}

The framework has four layers, depicted in
Figure~\ref{fig:architecture}. From inside out:

\begin{enumerate}
  \item \textbf{Inner particle filter.} For a fixed parameter vector
        $\thetadyn$, runs an SDE-driven bootstrap particle filter over
        the latent state. Returns an unbiased estimator of the marginal
        observation likelihood $p(y_{1:T} \mid \thetadyn)$.
        Implementation: \path{smc2fc/filtering/gk_dpf_v3_lite.py}.
  \item \textbf{Outer SMC$^2$ over parameters.} Runs adaptive tempered
        SMC over the parameter vector $\thetadyn$, calling the inner
        filter to evaluate likelihoods and using HMC moves between
        tempering steps to rejuvenate the particle cloud.
        Implementations: \path{smc2fc/core/tempered_smc.py} (BlackJAX
        wrapper, kept as a reference) and
        \path{smc2fc/core/jax_native_smc.py} (the production
        compile-once backend; see Section~\ref{sec:jax-native}).
  \item \textbf{Controller.} Reuses the outer SMC$^2$ engine, but with
        the data-likelihood replaced by a control-as-inference
        cost-likelihood $\exp(-\beta J(\thetactrl))$ over a finite-
        dimensional schedule parameter $\thetactrl$. Returns a
        posterior cloud over schedules; the posterior mean is the
        applied control. Implementation: \path{smc2fc/control/}.
  \item \textbf{MPC driver.} Orchestrates the receding-horizon loop:
        for each replan time, advance the plant under the previous
        stride's control, refresh the parameter posterior on the new
        observation window, then call the controller for the next
        stride. The driver is user code; the framework provides the
        ingredients.
\end{enumerate}

\begin{figure}[h]
\centering
\begin{tikzpicture}[
  node distance=8mm,
  every node/.style={font=\small},
  layer/.style={rectangle, rounded corners, draw, minimum height=8mm,
                minimum width=46mm, align=center},
  ext/.style={rectangle, draw, dashed, minimum height=8mm,
              minimum width=46mm, align=center, fill=gray!10},
  arrow/.style={-Stealth, thick},
]
\node[ext]   (model)   {User model: SDE + observation + prior};
\node[layer, below=of model] (pf) {Inner particle filter\\\path{smc2fc/filtering/}};
\node[layer, below=of pf]    (smc2) {Outer SMC$^2$ (HMC moves, adaptive tempering)\\\path{smc2fc/core/}};
\node[layer, below=of smc2]  (ctrl) {Controller (control-as-inference)\\\path{smc2fc/control/}};
\node[ext, below=of ctrl]    (mpc)  {User MPC driver: plant $\leftrightarrow$ filter $\leftrightarrow$ controller};

\draw[arrow] (model) -- node[right] {drift, obs, prior} (pf);
\draw[arrow] (pf)    -- node[right] {$\Lhat_N(\thetadyn)$ (unbiased)} (smc2);
\draw[arrow] (smc2)  -- node[right] {posterior cloud over $\thetadyn$} (ctrl);
\draw[arrow] (ctrl)  -- node[right] {posterior cloud over $\thetactrl$} (mpc);
\draw[arrow, dashed] (mpc.east) to[out=0, in=0, looseness=2] node[right] {observations} (smc2.east);
\end{tikzpicture}
\caption{The four layers of the \texttt{smc2fc} framework. Solid blocks
are framework code; dashed blocks are user code. The dashed arrow on
the right closes the receding-horizon loop: each stride's observations
become the next stride's likelihood input.}
\label{fig:architecture}
\end{figure}

\subsection{Reading order}

The document is structured to be readable straight through, but two
shortcuts are useful.

\paragraph{Mathematician's path.} If you want the mathematical anchors
first, read in this order:
Section~\ref{sec:setting} (the HMM setup),
Section~\ref{sec:pf} (Bayesian filtering recursion and the unbiased
likelihood property), Section~\ref{sec:smc2} (the tempered posterior),
Section~\ref{sec:control} (control-as-inference),
Section~\ref{sec:mpc} (receding-horizon principle).

\paragraph{Engineer's path.} If you want to use or modify the code,
read in this order:
Section~\ref{sec:setting} (the two interfaces you must conform to),
Section~\ref{sec:jax-native} (the production compile-once backend and
how to switch between it and the BlackJAX wrapper),
Section~\ref{sec:config} (every knob in \texttt{SMCConfig}, what each
one does and how to tune it), Section~\ref{sec:control} (the two
chance-constrained cost variants and which one to use),
Appendix~\ref{app:code-api} (one-line API lookup).

\subsection{Conventions}

\begin{itemize}
  \item $\thetadyn$ denotes the dynamics-model parameter vector
        (inferred from data); $\thetactrl$ denotes the schedule
        parameter vector (chosen by the controller). Both live in
        $\R^d$ for some user-defined $d$.
  \item $X_n$ is the latent state at step $n$, $Y_n$ the observation,
        $\Phi_n$ (or generic $U_n$) the exogenous control input at the
        same step.
  \item $N$ is the number of \emph{state} particles (inner PF);
        $M$ is the number of \emph{parameter} particles
        (outer SMC$^2$). Code uses \texttt{n\_pf\_particles} and
        \texttt{n\_smc\_particles} respectively.
  \item File:line citations are checked against the current source at
        the time of writing. If the line numbers drift (e.g.\ after a
        refactor), the function name remains the canonical reference.
\end{itemize}

\subsection{What this document does not promise}

The framework is research code, not a textbook implementation. It has
been tuned against a small number of model classes (the FSA
physiological models, the SWAT hydrological model). Hyperparameter
defaults are reasonable starting points for similar classes; for
genuinely different problems they will need re-tuning, and
Section~\ref{sec:config} is the diagnostic guide for doing so. Where
the framework has hit empirical limits, those limits are documented in
Sections~\ref{sec:jax-native}--\ref{sec:control} and
Section~\ref{sec:future-work}, not glossed.
